\chapter{主要符号表}

\begin{itemize}[align=left]
    \item[$\bm r,\bm r'$] 场点与源点的位置矢量。$x,y,z$ 是直角坐标分量，$\uvec r$ 是径向单位矢量。
    \item[$\bm E,\bm B$] 电场强度与磁感应强度。介质中的宏观场还用到电位移 $\bm D$、磁场强度 $\bm H$、极化强度 $\bm P$ 和磁化强度 $\bm M$。
    \item[$\Phi,\bm A,\emf$] 电势、磁矢势和电动势。计算磁通量时用 $\varPhi=\int_S\bm B\cdot\d\bm A$，其中 $\d\bm A$ 是有向面积元。
    \item[$\varepsilon_\relative,\mu_\relative$] 相对介电常数和相对磁导率。$\chi_\electric$ 是电极化率，$m_\electron$ 是电子质量。
    \item[$\theta_\incident,\theta_\reflected,\theta_\transmitted$] 入射角、反射角和折射角。界面问题中的下标 $\incident,\reflected,\transmitted$ 分别标记三束波。
    \item[$\theta_\critical,k_\cutoff,\bm E_\conductor$] 全反射临界角、波导截止波数和导体内的电场。波导中的 $\bm E_\transv$ 表示垂直于传播轴的场分量。
\end{itemize}
